\section{الفصل~\ref{sec:dishbrain}. \foreignlanguage{english}{DishBrain}}

بنية المنصة ونتائج اللعب مأخوذة من دراسة تجريبية واحدة ومن متابعتها؛ وصيغ الضجيج والانجراف نحو الضبوط الجيدة مشتقة في الفصل ومختبرة بنموذج الكتاب؛ أما آلية التعلم في الطبق فلم تُحدَّد.

{\footnotesize
\setlength{\tabcolsep}{3pt}\renewcommand{\arraystretch}{1.15}
\begin{longtable}{>{\raggedright}p{0.5cm}>{\raggedright}p{4.0cm}>{\raggedright}p{5.6cm}>{\raggedright}p{2.3cm}>{\raggedright\arraybackslash}p{1.7cm}}
\toprule
م & القضية & التجربة وما قيس & المصدر & الحالة \\
\midrule\endfirsthead
\toprule
م & القضية & التجربة وما قيس & المصدر & الحالة \\
\midrule\endhead
1 & المنصة: نحو $800\,000$ خلية من قشرة الفأر أو نحو $10^6$ خلية بشرية على الشريحة؛ $26\,000$ قطب على $8\,\text{mm}^2$، وتسجيل حتى $1024$، وتحفيز عبر $8$ & طرائق الدراسة: شريحة \foreignlanguage{english}{MaxOne}، $26\,000$ قطب من البلاتين على $8\,\text{mm}^2$، وتسجيل $1024$ قناة بـ$20\,000$ عينة في الثانية؛ ويمكن تقنيًا تحفيز حتى $32$ قطبًا، ونجح التحفيز المستقل عبر $8$؛ واستُعملت مزارع الفأر من نحو $10$ أيام & \cite{Kagan2022} & مُقاس \\
2 & حلقة بنبضة ساعة $10\ms$: نبضات المنطقتين الحركيتين تحرّك المضرب (الشكل~\ref{fig:dishloop}) & تُعدّ النبضات في نوافذ مدة كل منها $10\ms$؛ والتأخير من النبضة إلى المنبه المجيب نحو $5\ms$، يحدده مخزن العتاد المؤقت & \cite{Kagan2022} & مُقاس \\
3 & الدخل: ترميز المكان (أيّ الأقطاب الثمانية) وترميز التردد $4\text{--}40$ نبضة/ث & في التجربة الرئيسية يزيد تردد التحفيز خطيًا من $4\,\text{Hz}$ حين تكون الكرة عند الجدار البعيد إلى $40\,\text{Hz}$ عند المضرب؛ وسعة نبضات الكرة $75\mV$؛ والنبضات مستطيلة ثنائية الطور & \cite{Kagan2022} & مُقاس \\
4 & الخرج $\Delta p=c\,(n_\uparrow-n_\downarrow)$~\eqref{eq:dishreadout}، وضجيج العدّ في النافذة $1/\sqrt{RT}$ & قاعدة الفرق بين المنطقتين موصوفة في الدراسة (بأوزان للمنطقتين بحسب الفعالية)، والصيغة الدقيقة غير معطاة. وتقدير الضجيج نتيجة للعدّ البواسوني، مشتق في الفصل & \cite{Kagan2022}؛ استنتاج في الفصل & نظرية \\
5 & المعلّم: بعد الإصابة منبه متوقَّع $75\mV$، $100$ نبضة/ث، $0.1\,\text{s}$؛ وبعد الإخفاق منبه غير متوقَّع، $150\mV$، $5$ نبضات/ث، $4\,\text{s}$ في أماكن ولحظات عشوائية، ثم $4\,\text{s}$ توقف & الطرائق: بعد الإصابة $75\mV$، $100\,\text{Hz}$، $100\ms$ على الأقطاب الثمانية كلها معًا؛ وبعد الإخفاق $150\mV$، $5\,\text{Hz}$ على أقطاب عشوائية في لحظات عشوائية مدة $4\,\text{s}$، ثم $4\,\text{s}$ بلا تحفيز. ولا دوبامين في المزرعة & \cite{Kagan2022} & مُقاس \\
6 & الشبكة تتصرف بحيث يصير دخلها متوقَّعًا & مبدأ الطاقة الحرة اقتراح نظري اشتق منه الباحثون البروتوكول؛ والتجربة تتفق معه، لكنها لا تميّزه من آليات أبسط (السطران~\babelsublr{7--8}) & \cite{Friston2010,Kagan2022} & فرضية \\
7 & إذا كان الفشل يهزّ الشبكة والنجاح لا يهزّها، فالوقت في الضبط $\rho\propto1/P_{\text{miss}}$~\eqref{eq:dishdensity} (القضية~\ref{prop:dishscramble}) & ينتج من توازن التدفقات في سلسلة ماركوف بدفعات متناظرة، مشتق في الفصل. ولا اختبار مباشر على مزرعة & استنتاج في الفصل & نظرية \\
8 & نموذج ”الخلط عند الإخفاق“ يتعلم: نسبة الإصابات $0.21\to0.38$؛ ومع دفعة بعد كل جولة $\approx0.12$، وبلا دفعات $0.19$ (الشكل~\ref{fig:dishmodel}) & حساب، النص البرمجي \foreignlanguage{english}{\texttt{models/dishbrain\_model.py}}، $40$ مزرعة نموذجية بـ$2000$ جولة لكل منها: $0.21\to0.38$؛ $0.13\to0.11$؛ $0.19\to0.19$ & --- & نموذج \\
9 & سلاسل الكرات المصدودة تطول في المزارع الحية ذات الحلقة الكاملة فقط؛ والضوابط لا تتعلم & $399$ جلسة مدة كل منها $20$ دقيقة: $9$ مزارع من الفأر ($101$ جلسة)، و$11$ بشرية ($138$)، و$6$ شرائح بوسط الاستنبات ($80$)، و$20$ مزرعة في الراحة ($42$)، و$3$ محاكيات ($38$). متوسط طول الجولة في آخر $15$ دقيقة أكبر منه في أول $5$ في المزارع الحية فقط؛ ولا أثر في الخلايا غير العصبونية (\foreignlanguage{english}{HEK293T}) ولا في الشرائح بوسط الاستنبات & \cite{Kagan2022} & مُقاس \\
10 & النمو ذو الدلالة داخل الجلسة مع التغذية الراجعة الكاملة فقط؛ ومع الصمت بدل المنبه يكون اللعب في آخر الجلسة أفضل مما بلا تغذية راجعة، لكن الأثر أصغر & $486$ جلسة على مدى $3$ أيام: ”المنبه“ ($n=27$)، و”الصمت“ ($17$)، و”بلا تغذية راجعة“ ($15$). النمو مع الزمن ذو دلالة في شرط ”المنبه“ فقط؛ وفي آخر الجلسة تفوّق شرط ”الصمت“ على ”بلا تغذية راجعة“ بأثر أصغر & \cite{Kagan2022} & مُقاس \\
11 & الأثر صغير؛ وهل تتعلم الشبكة تعلمًا دائمًا مسألة مفتوحة & داخل الجلسة التحسن متين، لكن الباحثين أنفسهم يقولون إن التعلم بين الجلسات على مدى أيام لم يُلاحَظ ثابتًا & \cite{Kagan2022} & مُقاس \\
12 & في أثناء اللعب تقترب الشبكة من النظام الحرج، وكلما اقتربت كان اللعب أفضل & تحليل الانهيارات المتسلسلة للفعالية في المزارع نفسها: الدخل المنظَّم يقرّب الشبكة من النظام الحرج، وجودة اللعب ترتبط بالقرب منه؛ لكن الحرجية وحدها، بلا معلومات عن عواقب الأفعال، لا تكفي للتعلم & \cite{Habibollahi2023} & مُقاس \\
13 & ”وعي“ المزرعة لم يُبيَّن & مقالة ردّ لثلاثين باحثًا: تغيّر السلوك في حلقة مغلقة لا يثبت ذكاءً ولا إحساسًا؛ ولا تجربة تبيّن ذلك & \cite{Balci2023} & فرضية \\
14 & التجربة الضابطة الرابعة المقترحة: منبه متوقَّع بعد الإخفاق بالمدة والطاقة نفسيهما (القسم~\ref{sec:dishspec}) & لم تُجرَ تجربة كهذه؛ هذا اقتراح الكتاب & --- & فرضية \\
\bottomrule
\end{longtable}}
